% Format buku TA FIF. Kompilasi menggunakan XeLaTeX.
\usepackage{fontspec}
\IfFontExistsTF{Aptos}{\setmainfont{Aptos}}{%
  \PackageError{buku-ta}{Font Aptos belum terpasang}{Pasang Microsoft Aptos Fonts sesuai README.md.}}
\newfontfamily\headingfont{Aptos Display}
\IfFontExistsTF{Times New Roman}{\newfontfamily\referencefont{Times New Roman}}{\newfontfamily\referencefont{TeX Gyre Termes}}
\usepackage{babel}
\babelprovide[import,main,hyphenrules=nohyphenation]{indonesian}
\usepackage[a4paper,left=4cm,right=3cm,top=3cm,bottom=3cm,
  headheight=15pt,headsep=38.81pt,footskip=.5cm]{geometry}
\usepackage{graphicx,amsmath,amssymb,array,tabularx,longtable,booktabs}
\usepackage[table]{xcolor}
\usepackage{float}
\usepackage[normalem]{ulem}
\usepackage{enumitem}
\usepackage{titlesec,titletoc,caption,fancyhdr}
\usepackage[numbers,square,sort&compress]{natbib}
\usepackage{bibentry}
\usepackage[hidelinks,hypertexnames=false,unicode]{hyperref}
\usepackage{bookmark}
\graphicspath{{gambar/}}
% Word: 12 pt, 1.5 lines = sekitar 21.95 pt pada Aptos; paragraph after 6 pt.
\makeatletter
\renewcommand\normalsize{\@setfontsize\normalsize{12}{21.95}%
  \abovedisplayskip 12pt plus 3pt minus 3pt
  \belowdisplayskip \abovedisplayskip}
\makeatother
\normalsize
\setlength{\parindent}{0pt}
\setlength{\parskip}{6pt}
\tolerance=9999
\setlength{\emergencystretch}{3em}
\setcounter{secnumdepth}{2}
\setcounter{tocdepth}{2}
\titleformat{\chapter}[block]
  {\headingfont\fontsize{14}{17.1}\selectfont\bfseries\centering}
  {BAB \thechapter}{.5em}{}
\titlespacing*{\chapter}{0pt}{-21.95pt}{30pt}
\titleformat{\section}{\normalsize\bfseries}{\thesection.}{.6em}{}
\titleformat{\subsection}{\normalsize\bfseries}{\thesubsection.}{.6em}{}
\titlespacing*{\section}{0pt}{6pt}{0pt}
\titlespacing*{\subsection}{0pt}{6pt}{0pt}
\renewcommand{\thefigure}{\arabic{chapter}.\arabic{figure}}
\renewcommand{\thetable}{\arabic{chapter}.\arabic{table}}
\renewcommand{\theequation}{\arabic{chapter}.\arabic{equation}}
\DeclareCaptionFont{templatecaption}{\fontsize{9}{11}\selectfont\bfseries}
\captionsetup{font=templatecaption,labelsep=period,justification=centering,
  singlelinecheck=false,skip=6pt}
% Aturan tertulis tabel: hanya nomenklatur bold, dua spasi, judul sentence case.
\DeclareCaptionFont{ukurantabel}{\fontsize{9}{11}\selectfont}
\DeclareCaptionLabelSeparator{duaspasi}{\hspace{\dimexpr2\fontdimen2\font\relax}}
\captionsetup[table]{position=top,font=ukurantabel,labelfont=bf,
  textfont=normalfont,labelsep=duaspasi}
\setlist{topsep=0pt,itemsep=6pt,parsep=0pt,partopsep=0pt,leftmargin=1.27cm}
\setlist[itemize,1]{label=\textbullet}
\setlist[itemize,2]{label=\textopenbullet}
\setlist[enumerate,1]{label=\alph*.}
\setlength{\tabcolsep}{5pt}
\renewcommand{\arraystretch}{1.2}
% Header tabel isi: latar hitam dan teks putih sesuai template Word.
\newcommand{\HeaderTabel}[1]{\textcolor{white}{\textbf{#1}}}
\fancypagestyle{awal}{\fancyhf{}\fancyfoot[C]{\fontsize{12}{14}\selectfont\thepage}%
  \renewcommand{\headrulewidth}{0pt}\renewcommand{\footrulewidth}{0pt}}
\fancypagestyle{isi}{\fancyhf{}\fancyhead[R]{\fontsize{12}{14}\selectfont\thepage}%
  \renewcommand{\headrulewidth}{0pt}\renewcommand{\footrulewidth}{0pt}}
\fancypagestyle{plain}{\fancyhf{}\fancyfoot[C]{\fontsize{12}{14}\selectfont\thepage}%
  \renewcommand{\headrulewidth}{0pt}\renewcommand{\footrulewidth}{0pt}}
\newcommand{\HalamanAwal}[1]{\chapter*{#1}\addcontentsline{toc}{chapter}{#1}}
\newcommand{\AwalIsi}{\clearpage\pagenumbering{arabic}\pagestyle{isi}}
% Daftar otomatis, tanpa menyalin nomor halaman statis dari Word.
\titlecontents{chapter}[0pt]{\fontsize{11}{19.45}\selectfont\bfseries}
  {BAB \thecontentslabel\quad}{}{\titlerule*[.7pc]{.}\contentspage}
\titlecontents{section}[1.3cm]{\fontsize{10}{18.2}\selectfont\bfseries}
  {\contentslabel{1.3cm}}{}{\titlerule*[.7pc]{.}\contentspage}
\titlecontents{subsection}[2.0cm]{\fontsize{10}{18.2}\selectfont\bfseries}
  {\contentslabel{1.4cm}}{}{\titlerule*[.7pc]{.}\contentspage}
\titlecontents{figure}[2.3cm]{\fontsize{11}{19.45}\selectfont}
  {\contentslabel[Gambar \thecontentslabel.]{2.3cm}}{}
  {\titlerule*[.7pc]{.}\contentspage}
\titlecontents{table}[2.0cm]{\fontsize{11}{19.45}\selectfont}
  {\contentslabel[Tabel \thecontentslabel.]{2.0cm}}{}
  {\titlerule*[.7pc]{.}\contentspage}
\addto\captionsindonesian{%
  \renewcommand{\contentsname}{DAFTAR ISI}%
  \renewcommand{\listfigurename}{DAFTAR GAMBAR}%
  \renewcommand{\listtablename}{DAFTAR TABEL}%
  \renewcommand{\bibname}{DAFTAR PUSTAKA}}

% Catatan kaki asli: nomor otomatis dan isi 10 pt pada bagian bawah halaman.
\makeatletter
\renewcommand\footnotesize{\@setfontsize\footnotesize{10}{12}}
\makeatother

% Contoh entri lengkap berbagi database BibTeX; anchor hanya dibuat di daftar pustaka.
\makeatletter
\newcommand{\ContohPustaka}[1]{%
  \cite{#1}\quad\begingroup
  \let\hyper@natanchorstart\@gobble
  \let\hyper@natanchorend\@empty
  \bibentry{#1}\endgroup}
\makeatother

\makeatletter
\long\def\@makefntext#1{\parindent 1.8em\noindent
  \hb@xt@1.8em{\hss\@makefnmark}#1}
\makeatother
